\chapter*{Conclusion and Future Work}
\addcontentsline{toc}{chapter}{Conclusion and Future Work}

This thesis studies Correlation-Aware Reward Exploration (CARE) in reinforcement learning. Instead of treating CARE as a method used with only one type of curiosity reward, this work views it as an adaptive layer to study how state-dependent scaling of intrinsic rewards affects learning. The thesis considers three representative intrinsic reward families: count-based bonuses, the Intrinsic Curiosity Module (ICM), and Rewarding Impact-Driven Exploration (RIDE). CARE is used as the mechanism that decides how much these rewards should influence learning in each state. By scaling intrinsic reward according to its within-batch correlation with the extrinsic GAE advantage of the value function, CARE offers a simple way to balance exploration and exploitation during training without using costly meta-gradient methods.

Experiments on six MiniGrid environments evaluate each intrinsic reward module under a sweep of fixed scaling coefficients $\beta \in \{0.0005,\, 0.001,\, 0.005,\, 0.01,\, 0.05\}$ and under the CARE adaptive variant, with PPO without intrinsic reward as a control, for a total of $570$ training runs at $10^6$ steps each. Across the eighteen (environment, module) cells, CARE matches or exceeds the best post-hoc fixed coefficient in $7$ cells (most notably CARE-RIDE on UnlockPickup at $0.536$ vs $0.369$ for the best fixed value, and CARE-COUNT on KeyCorridor at $0.831$ vs $0.828$), ties within $\pm 0.005$ raw reward in $4$ further cells, and underperforms in the remaining $7$. CARE-COUNT also achieves the lowest variance and the fastest sample-efficiency on Empty-16x16 ($66\text{k}$ steps to threshold vs $82\text{k}$ for the best fixed value). The benefits concentrate in the count-based and RIDE pairings; ICM is the exception, where CARE inherits the underlying signal's poor alignment with extrinsic learning slack and even collapses to zero on KeyCorridor. The largest fixed coefficient $\beta = 0.05$ in the sweep is universally catastrophic, which empirically justifies clipping $\beta_\text{max} = 5 \times 10^{-2}$ in the Beta Network. A second limitation appears on Empty-16x16: when extrinsic rewards are almost always missing, the GAE advantage carries no signal-strength information across states, the correlation gradient is approximately zero-mean noise, and the scaling factor falls back to its reference value $\beta_0$, behaving like a fixed coefficient. This suggests that the effectiveness of CARE depends both on the alignment between the intrinsic generator and the extrinsic-success path, and on how often the agent encounters extrinsic feedback that produces a non-degenerate advantage signal.

There are several directions for future work. First, the method can be applied to broader settings such as continuous control, multi-task RL, and embodied navigation, where the structure of intrinsic motivation differs from the gridworld setting studied here. Second, a clearer theoretical analysis of the correlation objective could help explain when adaptive scaling is most useful across different intrinsic reward types. Third, alternative supervision signals for the Beta Network beyond extrinsic GAE advantage---for example bootstrapped novelty or trajectory-level success indicators---may improve robustness in extremely sparse regimes where the correlation signal degenerates. These directions can further improve the use of adaptive intrinsic rewards in reinforcement learning.